\section{Quy trình nạp và kiểm soát chất lượng}
\subsection{Nạp raw và chuẩn hóa sang staging}
\code{etl_pipeline.py} đảm bảo database tồn tại, tạo schema \code{staging}, rồi COPY bảy tệp CSV vào \code{staging.raw_*} (HEADER, UTF8; toàn TEXT để tách việc đọc dữ liệu khỏi ép kiểu). Chế độ mặc định tạo lại staging và DW; tùy chọn \code{--reload-only} giữ schema và truncate rồi nạp lại, không phải nạp tăng dần.

\code{03_promote.sql} chuyển raw sang \code{staging.stg_*} bằng cách btrim định danh và chuỗi, dùng NULLIF để đổi chuỗi rỗng thành NULL, rồi ép kiểu timestamp, integer, numeric. Bước này không khử trùng lặp, không nội suy giá trị thiếu và không sửa số tiền; PK/FK và kiểu sẽ làm INSERT thất bại nếu dữ liệu không phù hợp. Pipeline in số dòng từng bảng staging để phát hiện mất dòng.

Thông tin giao hàng, kích thước sản phẩm, vị trí seller, loại thanh toán và số kỳ trả góp được bỏ khỏi staging có kiểu theo nhu cầu phân tích hiện có; do đó DW không hỗ trợ SLA giao hàng hay phân tích loại thanh toán.

\subsection{Nạp dimension trước fact}
\begin{enumerate}[leftmargin=*]
\item \code{dim_date}: sinh lịch liên tục từ ngày nhỏ nhất đến lớn nhất của ngày mua và ngày duyệt (mọi trạng thái); thứ và cuối tuần tính theo ISO.
\item \code{dim_customer}: DISTINCT theo \code{customer_unique_id}; \code{dim_product}: LEFT JOIN bản dịch danh mục; \code{dim_seller}: theo \code{seller_id}.
\item \code{fact_customer_order}: gom thanh toán theo \code{order_id} trước rồi LEFT JOIN; đơn không có payment giữ NULL, không thay bằng 0.
\item \code{fact_sales}: nối dòng hàng với đơn, khách và các dimension; lấy ngày mua và khách từ cùng đơn.
\end{enumerate}

\begin{table}[H]\centering\small
\begin{tabularx}{\linewidth}{p{3.7cm}YY}
\toprule Đích & Nguồn chính & Quy tắc đáng chú ý\\\midrule
\code{dim_customer} & stg\_customers & DISTINCT customer\_unique\_id, không phải customer\_id.\\
\code{dim_product} & stg\_products, bản dịch & LEFT JOIN; thiếu bản dịch vẫn giữ sản phẩm.\\
\code{fact_customer_order} & stg\_orders, customers, payments & Tổng tiền gom trước khi nối; địa lý theo customer\_id của đơn.\\
\code{fact_sales} & stg\_order\_items và các dimension & Giữ grain dòng hàng; không nối trực tiếp payment.\\
\bottomrule\end{tabularx}
\caption{Ánh xạ staging sang DW}
\end{table}

\subsection{Nghiệm thu dữ liệu}
Bộ nghiệm thu nằm ở \code{sql/dw/03_validate.sql}, chạy ngay sau khi nạp; mỗi tiêu chí sai sẽ \code{RAISE EXCEPTION}, pipeline dừng và toàn bộ lần nạp bị rollback. Các tiêu chí kiểm tra bất biến của dữ liệu, không gõ cứng số đếm của một bản dữ liệu cụ thể:

\begin{table}[H]\centering\small
\begin{tabularx}{\linewidth}{lY}
\toprule Nhóm & Nội dung kiểm tra\\\midrule
F1 & Tổng tiền hàng và cước ở \code{fact_sales} bằng tổng ở staging.\\
F2 & Không trùng cặp (order\_id, order\_item\_id); mọi đơn của \code{fact_sales} có trong fact đơn; ngày mua và khách khớp giữa hai fact.\\
F3 & Khóa ngoại tới các dimension đều ánh xạ được.\\
F4 & Tổng thanh toán theo đơn khớp staging; đơn có thanh toán NULL thì nguồn thực sự không có dòng thanh toán (số đơn như vậy chỉ được báo, không bắt buộc bằng một giá trị).\\
F5 & Số đơn theo từng trạng thái khớp giữa staging và fact.\\
F6 & Các cột của \code{dim_date} nhất quán với \code{full_date}.\\
F7 & Không có giá hoặc cước âm; không có khóa bắt buộc bị NULL.\\
\bottomrule\end{tabularx}
\caption{Tiêu chí nghiệm thu F1--F7}
\end{table}

Có thể chạy lại bộ nghiệm thu mà không nạp lại dữ liệu bằng \code{check_acceptance.py}; script thoát mã 1 khi có tiêu chí sai. Các phép so sánh dùng \code{IS DISTINCT FROM} nên bắt được cả trường hợp một vế NULL. Một phép thử làm rỗng \code{fact_sales} trong giao dịch rồi hủy cho thấy bản dùng \code{<>} vẫn báo đạt (vì tổng tiền là NULL), còn bản hiện tại báo lỗi F1.

\subsection{Giao dịch và khả năng tái lập}
Toàn bộ pipeline (tạo schema, nạp raw, promote, tạo và nạp DW, nghiệm thu, tạo view DSS) chạy trong một giao dịch PostgreSQL và chỉ \code{COMMIT} khi mọi bước thành công. Lỗi ở bất kỳ bước nào làm giao dịch bị rollback, nên một lần chạy lại thất bại vẫn giữ nguyên kho cũ. Điều này đã được kiểm tra bằng cách giả lập lỗi ở bước nghiệm thu, khi đó pipeline dừng và \code{fact_sales} vẫn giữ 112.650 dòng, view DSS vẫn tồn tại.

Chế độ mặc định xóa schema rồi tạo lại nên cho cùng kết quả sau mỗi lần chạy; \code{--reload-only} không RESTART IDENTITY nên khóa thay thế có thể khác và không được dùng làm định danh nghiệp vụ. Vì tạo lại schema \code{dw} xóa cả các view, bước cuối của pipeline tạo lại view \code{dw.v_dss_monthly_demand} (mục~\ref{sec:dss-data}). Pipeline đã được chạy đầy đủ trên PostgreSQL 16 và đạt F1--F7 (chương~\ref{sec:results}).
